\documentclass[10pt,a4paper]{amsart}
\usepackage{amsmath,amssymb,amsthm,graphicx}
\usepackage{verbatim}
\newenvironment{reference}{\comment}{\endcomment}
%\newenvironment{reference}{}{}
\newenvironment{slogan}{\comment}{\endcomment}
\newenvironment{history}{\comment}{\endcomment}
% For commutative diagrams we use Xy-pic
\usepackage[all]{xy}

% We use 2cell for 2-commutative diagrams.
\xyoption{2cell}
% Theorem environments.
%
\theoremstyle{plain}
\newtheorem{theorem}[subsection]{Theorem}
\newtheorem{proposition}[subsection]{Proposition}
\newtheorem{lemma}[subsection]{Lemma}

\theoremstyle{definition}
\newtheorem{definition}[subsection]{Definition}
\newtheorem{example}[subsection]{Example}
\newtheorem{exercise}[subsection]{Exercise}
\newtheorem{situation}[subsection]{Situation}

\theoremstyle{remark}
\newtheorem{remark}[subsection]{Remark}
\newtheorem{remarks}[subsection]{Remarks}

\numberwithin{equation}{subsection}

% Macros
%
\def\lim{\mathop{\mathrm{lim}}\nolimits}
\def\colim{\mathop{\mathrm{colim}}\nolimits}
\def\Spec{\mathop{\mathrm{Spec}}}
\def\Hom{\mathop{\mathrm{Hom}}\nolimits}
\def\Ext{\mathop{\mathrm{Ext}}\nolimits}
\def\SheafHom{\mathop{\mathcal{H}\!\mathit{om}}\nolimits}
\def\SheafExt{\mathop{\mathcal{E}\!\mathit{xt}}\nolimits}
\def\Sch{\mathit{Sch}}
\def\Mor{\mathop{\mathrm{Mor}}\nolimits}
\def\Ob{\mathop{\mathrm{Ob}}\nolimits}
\def\Sh{\mathop{\mathit{Sh}}\nolimits}
\def\NL{\mathop{N\!L}\nolimits}
\def\CH{\mathop{\mathrm{CH}}\nolimits}
\def\proetale{{pro\text{-}\acute{e}tale}}
\def\etale{{\acute{e}tale}}
\def\QCoh{\mathit{QCoh}}
\def\Ker{\mathop{\mathrm{Ker}}}
\def\Im{\mathop{\mathrm{Im}}}
\def\Coker{\mathop{\mathrm{Coker}}}
\def\Coim{\mathop{\mathrm{Coim}}}

% Boxtimes
%
\DeclareMathSymbol{\boxtimes}{\mathbin}{AMSa}{"02}

%
% Macros for moduli stacks/spaces
%
\def\QCohstack{\mathcal{QC}\!\mathit{oh}}
\def\Cohstack{\mathcal{C}\!\mathit{oh}}
\def\Spacesstack{\mathcal{S}\!\mathit{paces}}
\def\Quotfunctor{\mathrm{Quot}}
\def\Hilbfunctor{\mathrm{Hilb}}
\def\Curvesstack{\mathcal{C}\!\mathit{urves}}
\def\Polarizedstack{\mathcal{P}\!\mathit{olarized}}
\def\Complexesstack{\mathcal{C}\!\mathit{omplexes}}
% \Pic is the operator that assigns to X its picard group, usage \Pic(X)
% \Picardstack_{X/B} denotes the Picard stack of X over B
% \Picardfunctor_{X/B} denotes the Picard functor of X over B
\def\Pic{\mathop{\mathrm{Pic}}\nolimits}
\def\Picardstack{\mathcal{P}\!\mathit{ic}}
\def\Picardfunctor{\mathrm{Pic}}
\def\Deformationcategory{\mathcal{D}\!\mathit{ef}}
\usepackage{cleveref}
\crefname{originallemma}{original lemma}{original lemmas}
\Crefname{originallemma}{Original lemma}{Original lemmas}
\makeatletter
\newlabel{lemma-composition-finite-type}{{lemma-composition-finite-type}{source}{}{sourcepointer-1}{}}
\newlabel{lemma-composition-finite-type@cref}{{[originallemma][0][]lemma-composition-finite-type}{[1][0][]source}{}{}{}}
\makeatother
\crefname{originalreference}{original reference}{original references}
\Crefname{originalreference}{Original reference}{Original references}
\makeatletter
\newlabel{morphisms-lemma-dimension-fibre-specialization}{{morphisms-lemma-dimension-fibre-specialization}{source}{}{sourcepointer-2}{}}
\newlabel{morphisms-lemma-dimension-fibre-specialization@cref}{{[originalreference][0][]morphisms-lemma-dimension-fibre-specialization}{[1][0][]source}{}{}{}}
\makeatother
\crefname{originalreference}{original reference}{original references}
\Crefname{originalreference}{Original reference}{Original references}
\makeatletter
\newlabel{morphisms-lemma-quasi-finite-at-point-characterize}{{morphisms-lemma-quasi-finite-at-point-characterize}{source}{}{sourcepointer-3}{}}
\newlabel{morphisms-lemma-quasi-finite-at-point-characterize@cref}{{[originalreference][0][]morphisms-lemma-quasi-finite-at-point-characterize}{[1][0][]source}{}{}{}}
\makeatother
\crefname{originalreference}{original reference}{original references}
\Crefname{originalreference}{Original reference}{Original references}
\makeatletter
\newlabel{spaces-morphisms-lemma-compare-tr-deg}{{spaces-morphisms-lemma-compare-tr-deg}{source}{}{sourcepointer-4}{}}
\newlabel{spaces-morphisms-lemma-compare-tr-deg@cref}{{[originalreference][0][]spaces-morphisms-lemma-compare-tr-deg}{[1][0][]source}{}{}{}}
\makeatother
\crefname{originalreference}{original reference}{original references}
\Crefname{originalreference}{Original reference}{Original references}
\makeatletter
\newlabel{stacks-properties-lemma-points-cartesian}{{stacks-properties-lemma-points-cartesian}{source}{}{sourcepointer-5}{}}
\newlabel{stacks-properties-lemma-points-cartesian@cref}{{[originalreference][0][]stacks-properties-lemma-points-cartesian}{[1][0][]source}{}{}{}}
\makeatother
\crefname{originalreference}{original reference}{original references}
\Crefname{originalreference}{Original reference}{Original references}
\makeatletter
\newlabel{stacks-properties-section-properties-morphisms}{{stacks-properties-section-properties-morphisms}{source}{}{sourcepointer-6}{}}
\newlabel{stacks-properties-section-properties-morphisms@cref}{{[originalreference][0][]stacks-properties-section-properties-morphisms}{[1][0][]source}{}{}{}}
\makeatother
\title{A stack locally quasi-finite over a field is a disjoint union of points with finite automorphism groups}
\author{The Stacks Project Authors}
\date{stacks-project:stacks-morphisms.tex:local-snapshot-20261012; GNU FDL 1.2 or later}
\begin{document}
\maketitle
\noindent\textit{Editorial note.} All delivered mathematical text is the author's own, copied from the pinned Stacks Project chapter source sources/stacks-algebraic-stacks-morphisms/pinned-original-source.tex; no statement or proof step was rewritten, reordered or removed.\par\medskip
\noindent\textit{Editorial note.} Preamble provenance: the Stacks Project keeps one single shared preamble for every chapter; this card takes it directly from the staged sources/stacks-algebraic-stacks-morphisms/preamble.tex as three verbatim spans (lines 8-12, 14-18 and 154-217), so no derived file was written and nothing was added or redefined. The wrapper class and the base packages are supplied by the assembler.\par\medskip
\noindent\textit{Editorial note.} Bibliography provenance: neither the counted statement nor its proof contains a citation (this lane rescans the whole span before writing the card), so the card carries no bibliography and no entry was reconstructed.\par\medskip
\noindent\textit{Editorial note.} Selection provenance: this card was opened under the file-based control contract (CONTROL.md), with the mechanical pre-card gate run on its statement and proof interval before the card was written.\par\medskip
\noindent\textit{Allowed standard prerequisites (editorial).}\begin{itemize}
\item The characterization of local quasi-finiteness through locally finite type and discrete field fibres. Invoked inside the counted proof (source lines 5010-5018). Gives discreteness of the space from the hypothesis.
\item The description of a locally finite type stack over a field through a presentation by a group action. Used inside the counted proof (source lines 5019-5030). Exhibits the stack as a quotient of points by finite groups.
\end{itemize}\medskip
\ifdefined\hypertarget\hypertarget{sourcepointer-1}{}\fi
\noindent\textit{Source reference.} Original lemma lemma-composition-finite-type: sources/stacks-algebraic-stacks-morphisms/pinned-original-source.tex, source lines 3488--3488. Navigation pointer only: the counted text refers to the Stacks item lemma-composition-finite-type, which is an existing labeled result of this same chapter and is not part of this card. The pointer carries the target's own label as its number and its definition line as its span; its content is not reproduced here.\par\medskip
\ifdefined\hypertarget\hypertarget{sourcepointer-2}{}\fi
\noindent\textit{Source reference.} Original reference morphisms-lemma-dimension-fibre-specialization: sources/stacks-algebraic-stacks-morphisms/pinned-original-source.tex, source lines 5018--5018. Navigation pointer only: the counted text refers to the Stacks item morphisms-lemma-dimension-fibre-specialization, which belongs to another chapter of the project and is not part of this card. The pointer carries the target's own label as its number because this chapter's staged snapshot does not contain the other chapter's tag table; its content is not reproduced here.\par\medskip
\ifdefined\hypertarget\hypertarget{sourcepointer-3}{}\fi
\noindent\textit{Source reference.} Original reference morphisms-lemma-quasi-finite-at-point-characterize: sources/stacks-algebraic-stacks-morphisms/pinned-original-source.tex, source lines 5016--5016. Navigation pointer only: the counted text refers to the Stacks item morphisms-lemma-quasi-finite-at-point-characterize, which belongs to another chapter of the project and is not part of this card. The pointer carries the target's own label as its number because this chapter's staged snapshot does not contain the other chapter's tag table; its content is not reproduced here.\par\medskip
\ifdefined\hypertarget\hypertarget{sourcepointer-4}{}\fi
\noindent\textit{Source reference.} Original reference spaces-morphisms-lemma-compare-tr-deg: sources/stacks-algebraic-stacks-morphisms/pinned-original-source.tex, source lines 5030--5030. Navigation pointer only: the counted text refers to the Stacks item spaces-morphisms-lemma-compare-tr-deg, which belongs to another chapter of the project and is not part of this card. The pointer carries the target's own label as its number because this chapter's staged snapshot does not contain the other chapter's tag table; its content is not reproduced here.\par\medskip
\ifdefined\hypertarget\hypertarget{sourcepointer-5}{}\fi
\noindent\textit{Source reference.} Original reference stacks-properties-lemma-points-cartesian: sources/stacks-algebraic-stacks-morphisms/pinned-original-source.tex, source lines 2126--2126. Navigation pointer only: the counted text refers to the Stacks item stacks-properties-lemma-points-cartesian, which belongs to another chapter of the project and is not part of this card. The pointer carries the target's own label as its number because this chapter's staged snapshot does not contain the other chapter's tag table; its content is not reproduced here.\par\medskip
\ifdefined\hypertarget\hypertarget{sourcepointer-6}{}\fi
\noindent\textit{Source reference.} Original reference stacks-properties-section-properties-morphisms: sources/stacks-algebraic-stacks-morphisms/pinned-original-source.tex, source lines 248--248. Navigation pointer only: the counted text refers to the Stacks item stacks-properties-section-properties-morphisms, which belongs to another chapter of the project and is not part of this card. The pointer carries the target's own label as its number because this chapter's staged snapshot does not contain the other chapter's tag table; its content is not reproduced here.\par\medskip
\begin{lemma}
\label{lemma-locally-quasi-finite-over-field}
Let $\mathcal{X} \to \Spec(k)$ be a locally quasi-finite morphism
where $\mathcal{X}$ is an algebraic stack and $k$ is a field.
Let $f : V \to \mathcal{X}$ be a locally quasi-finite morphism where
$V$ is a scheme. Then $V \to \Spec(k)$ is locally quasi-finite.
\end{lemma}

\begin{proof}
By
Lemma \ref{lemma-composition-finite-type}
we see that $V \to \Spec(k)$ is locally of finite type.
Assume, to get a contradiction, that $V \to \Spec(k)$ is not
locally quasi-finite. Then there exists a nontrivial specialization
$v \leadsto v'$ of points of $V$, see
Morphisms, Lemma \ref{morphisms-lemma-quasi-finite-at-point-characterize}.
In particular $\text{trdeg}_k(\kappa(v)) > \text{trdeg}_k(\kappa(v'))$, see
Morphisms, Lemma \ref{morphisms-lemma-dimension-fibre-specialization}.
Because $|\mathcal{X}|$ is discrete we see that $|f|(v) = |f|(v')$.
Consider $R = V \times_\mathcal{X} V$. Then $R$ is an algebraic space
and the projections $s, t : R \to V$ are locally quasi-finite as base
changes of $V \to \mathcal{X}$ (which is representable by algebraic spaces
so this follows from the discussion in
Properties of Stacks, Section
\ref{stacks-properties-section-properties-morphisms}).
By
Properties of Stacks, Lemma \ref{stacks-properties-lemma-points-cartesian}
we see that there exists an $r \in |R|$ such that $s(r) = v$ and $t(r) = v'$.
By
Morphisms of Spaces, Lemma \ref{spaces-morphisms-lemma-compare-tr-deg}
we see that the transcendence degree of $v/k$ is equal to the
transcendence degree of $r/k$ is equal to the transcendence degree of
$v'/k$. This contradiction proves the lemma.
\end{proof}

\end{document}
